\chapter{Task 3: Fixed-Mesh Baseline Verification and Reference Reproduction}
\label{ch:baseline}

\section{Purpose of the baseline verification}
Before introducing mesh adaptation, the phase-field implementation must be rigorously verified against reference problems with established analytical or published numerical solutions. This chapter establishes source-code verification and reproduces the reference fracture response without mesh refinement. The staggered UEL/UMAT formulation of Moln\'ar and Gravouil \citep{Molnar2017} provides the implementation foundation, while the finite-size single-edge-notched Mode-I tensile benchmark of Pandey and Kumar \citep{Pandey2025} provides the authoritative quantitative comparison for the present research.

\section{Moln\'ar one-element constitutive verification}
The unchanged single-element test case supplied with the Moln\'ar--Gravouil implementation was compiled and executed to verify basic user-subroutine mechanics before running full fracture benchmarks. The calculation executed cleanly with $1{,}001$ increments in the step \texttt{Static}. Figure~\ref{fig:molnar_one_element} displays direct Abaqus CAE exports from the archived ODB, demonstrating that user subroutine \texttt{f42\_mixed\_uel.for} compiles, links, solves, and correctly outputs the phase-field variable \texttt{SDV15}.

\begin{figure}[htbp]
\centering
\begin{minipage}[t]{0.48\textwidth}
  \centering\includegraphics[width=\linewidth]{fig02a_molnar_one_element_mesh.png}\\[-0.3em]
  \small (a) Single-element physical discretization
\end{minipage}\hfill
\begin{minipage}[t]{0.48\textwidth}
  \centering\includegraphics[width=\linewidth]{fig02a_molnar_one_element_sdv15.png}\\[-0.3em]
  \small (b) Final \texttt{SDV15} field ($0 \le d \le 1$)
\end{minipage}
\caption{Direct Abaqus CAE output from the Moln\'ar single-element test ODB, verifying source compilation, element execution, and state variable reporting.}
\label{fig:molnar_one_element}
\end{figure}

\section{Moln\'ar--Gravouil single-edge-notched reproduction}
The author-supplied single-edge-notched tension specimen was reproduced next. Figure~\ref{fig:molnar_single_notch_mesh} shows the finite element mesh parsed directly from the original input deck, featuring a uniformly refined horizontal crack band, initial notch, and prescribed boundary conditions.

\begin{figure}[htbp]
\centering
\includegraphics[width=0.78\textwidth]{fig02b_molnar_single_notch_mesh_bc.png}
\caption{Finite element mesh and boundary conditions for the author-supplied Moln\'ar single-edge-notched specimen, reconstructed directly from the input deck.}
\label{fig:molnar_single_notch_mesh}
\end{figure}

The phase-field evolution across four representative displacement levels is illustrated in Figure~\ref{fig:molnar_phase_evolution}. The plotted field is \texttt{SDV15} ($d \in [0, 1]$) rendered in Abaqus/Viewer using a consistent viewport and fixed display range. As the applied vertical displacement increases, damage localizes smoothly from the notch tip ($u = 0.0050\,\mathrm{mm}$) into a fully developed crack ($d = 1.0$) propagating horizontally across the ligament ($u = 0.0060\text{--}0.0070\,\mathrm{mm}$).

\begin{figure}[htbp]
\centering
\begin{minipage}[t]{0.24\textwidth}\centering
\includegraphics[width=\linewidth]{fig02c_cae_molnar_u0020.png}\\[-0.3em]\small (a) $u=0.0020\,\mathrm{mm}$
\end{minipage}\hfill
\begin{minipage}[t]{0.24\textwidth}\centering
\includegraphics[width=\linewidth]{fig02c_cae_molnar_u0050.png}\\[-0.3em]\small (b) $u=0.0050\,\mathrm{mm}$
\end{minipage}\hfill
\begin{minipage}[t]{0.24\textwidth}\centering
\includegraphics[width=\linewidth]{fig02c_cae_molnar_u0060.png}\\[-0.3em]\small (c) $u=0.0060\,\mathrm{mm}$
\end{minipage}\hfill
\begin{minipage}[t]{0.24\textwidth}\centering
\includegraphics[width=\linewidth]{fig02c_cae_molnar_u0070.png}\\[-0.3em]\small (d) $u=0.0070\,\mathrm{mm}$
\end{minipage}
\caption{Direct Abaqus CAE phase-field contours (\texttt{SDV15}) at four converged displacement levels for the Moln\'ar single-notch specimen, plotted on the fixed physical range $d \in [0, 1]$.}
\label{fig:molnar_phase_evolution}
\end{figure}

Quantitative extraction of the reaction force versus applied displacement curve (Figure~\ref{fig:molnar_rf_u}) yielded a peak load of $F_{\max} = 0.727608\,\mathrm{kN}$ at $u = 0.006100\,\mathrm{mm}$, compared to the digitized reference peak of $0.690591\,\mathrm{kN}$ at $u = 0.005584\,\mathrm{mm}$ reported in \citep{Molnar2017} for $\ell_0 = 0.015\,\mathrm{mm}$. 

\begin{figure}[htbp]
\centering
\includegraphics[width=0.76\textwidth]{fig02d_molnar_single_notch_rf_u.png}
\caption{Reaction force versus prescribed displacement comparison between the author-supplied Moln\'ar model and the digitized published reference curve ($\ell_0 = 0.015\,\mathrm{mm}$).}
\label{fig:molnar_rf_u}
\end{figure}

To assess spatial convergence, a controlled mesh-refinement study was conducted across three meshes: coarse H0 ($3{,}930$ elements), intermediate H1 ($12{,}064$ elements), and fine H2 ($33{,}852$ elements), shown in Figure~\ref{fig:molnar_h_meshes}. The resulting force-displacement curves (Figure~\ref{fig:molnar_h_convergence}) confirmed monotonic convergence towards the continuum limit as element size decreased below $h \le \ell_0 / 2$.

\begin{figure}[htbp]
\centering
\begin{minipage}[t]{0.32\textwidth}\centering
\includegraphics[width=\linewidth]{fig02f_molnar_mesh_h0.png}\\[-0.3em]\small (a) H0 ($3{,}930$ elements)
\end{minipage}\hfill
\begin{minipage}[t]{0.32\textwidth}\centering
\includegraphics[width=\linewidth]{fig02f_molnar_mesh_h1.png}\\[-0.3em]\small (b) H1 ($12{,}064$ elements)
\end{minipage}\hfill
\begin{minipage}[t]{0.32\textwidth}\centering
\includegraphics[width=\linewidth]{fig02f_molnar_mesh_h2.png}\\[-0.3em]\small (c) H2 ($33{,}852$ elements)
\end{minipage}
\caption{Discretizations used in the Moln\'ar Mode-I mesh-convergence study.}
\label{fig:molnar_h_meshes}
\end{figure}

\begin{figure}[htbp]
\centering
\includegraphics[width=0.88\textwidth]{fig02e_molnar_h_convergence.png}
\caption{Force-displacement response across the H0/H1/H2 mesh sequence, illustrating convergence of peak force and softening trajectory.}
\label{fig:molnar_h_convergence}
\end{figure}

\section{Authoritative fixed-mesh Mode-I reproduction (Task 3)}
The principal benchmark used for evaluating mesh refinement is the Mode-I single-edge-crack square plate of Pandey \& Kumar \citep{Pandey2025}. The specimen geometry is a $1.0 \times 1.0\,\mathrm{mm}^2$ square domain containing an initial horizontal crack of length $a_0 = 0.5\,\mathrm{mm}$ located along the symmetry line $y = 0.5\,\mathrm{mm}$ (Figure~\ref{fig:mode1_geometry}). Material and fracture parameters are:
\begin{equation}
E = 210.0\,\mathrm{GPa}, \quad \nu = 0.30, \quad G_c = 2.7 \times 10^{-3}\,\mathrm{kN/mm} = 2700\,\mathrm{J/m^2}, \quad \ell_0 = 0.0075\,\mathrm{mm}, \quad \eta = 10^{-7}.
\end{equation}
The boundary conditions prescribe zero vertical displacement $u_y = 0$ along the bottom edge ($y = 0$), pin the bottom-right corner ($u_x = 0$ at $x=1, y=0$), and apply a uniform vertical tensile displacement $u_y = u$ to the top edge ($y = 1.0\,\mathrm{mm}$) via reference-point kinematic coupling. In the reference paper, the standard fixed mesh comprises a coarse global element size of $h_{\mathrm{cms}} = 0.02\,\mathrm{mm}$ and a uniformly refined horizontal crack corridor of height $0.1\,\mathrm{mm}$ with local element size $h = 0.003\,\mathrm{mm}$ ($15{,}192$ physical elements).

\begin{figure}[htbp]
\centering
\includegraphics[width=0.62\textwidth]{fig_mode1_geometry.pdf}
\caption{Geometry, loading, and boundary conditions for the Mode-I single-edge-crack tension benchmark of Pandey \& Kumar (2025).}
\label{fig:mode1_geometry}
\end{figure}

The fixed-mesh reference simulation was executed on the HPC cluster under Job \texttt{1398090.mmaster02}. The calculation converged smoothly through $7{,}000$ fixed increments with zero cutbacks (\texttt{Exit status 0}), yielding the results summarized in Table~\ref{tab:baseline_comparison}.

\begin{table}[htbp]
\centering
\caption{Quantitative validation of the fixed-mesh Mode-I reference reproduction (Job \texttt{1398090.mmaster02}) against published benchmark targets.}
\label{tab:baseline_comparison}
\begin{tabular}{lccc}
\toprule
\textbf{Metric} & \textbf{Published Target} \citep{Pandey2025} & \textbf{Reconstructed Baseline} & \textbf{Relative Difference} \\
\midrule
Peak Reaction Force $F_{\text{peak}}$ & $0.758000\,\mathrm{kN}$ & $0.757778\,\mathrm{kN}$ & $\mathbf{-0.029\%}$ \\
Displacement at Peak $u_{\text{peak}}$ & $0.005860\,\mathrm{mm}$ & $0.005857\,\mathrm{mm}$ & $\mathbf{-0.051\%}$ \\
Initial Elastic Stiffness $K_0$ & $\approx 138.0\,\mathrm{kN/mm}$ & $138.0837\,\mathrm{kN/mm}$ & $+0.06\%$ \\
Post-Peak Load Drop & $\approx 100\%$ & $99.97\%$ ($F_{\text{final}} = 0.232\,\mathrm{N}$) & Consistent \\
Physical Elements $N_{\mathrm{phys}}$ & $\approx 15{,}000$ & $15{,}192$ & $+1.28\%$ \\
Solver Increments & $7{,}000$ & $7{,}000$ & Exact \\
Solver Exit Status & Completed & $0$ (Normal) & Exact \\
\bottomrule
\end{tabular}
\end{table}

The relative errors of $-0.029\%$ in peak reaction force and $-0.051\%$ in peak displacement confirm an exceptional degree of quantitative agreement with the literature benchmark. This validates the underlying UEL implementation and provides the authoritative, frozen baseline for all subsequent adaptive remeshing evaluations (Task 3: \textbf{COMPLETE / SCIENTIFICALLY ACCEPTED}).

\section{Crack geometry verification: Sharp seam vs blunt notch}
During initial model development, the geometric representation of the initial crack was identified as a critical factor. An exploratory model that modeled the pre-crack as an explicit blunt slot of finite width $w = 0.001\,\mathrm{mm}$ exhibited a severely reduced initial elastic stiffness of $K_0 = 75.47\,\mathrm{kN/mm}$ and an artificial peak force of $0.1963\,\mathrm{kN}$ (Job \texttt{1398807.mmaster02}). 

Re-modeling the pre-crack as a true mathematical zero-gap seam using Abaqus \texttt{assignSeam} (with $100$ coincident node pairs along the crack flank and a unique crack-tip node \#145) restored the initial stiffness to $K_0 = 138.08\,\mathrm{kN/mm}$. This confirmed that the pre-crack must be modeled as a sharp crack seam to avoid artificial compliance and corner stress singularities.
